\chapter*{Введение}
\addcontentsline{toc}{chapter}{Введение}

Это личный конспект и собрание разборов доказательств на естественном языке к книге
Emily Riehl \emph{Category Theory in Context}. Конспект ведётся параллельно с формализацией
той же книги на языке Lean~4 в этом же репозитории (директория
\path{CategoryTheoryInContextLean/}) и служит двум целям.

Во-первых, это способ разобраться в материале книги медленнее и внимательнее, чем позволяет
Lean: записать доказательство обычными словами — честная проверка того, понимаю ли я, почему
утверждение верно. Тактика способна закрыть цель и без такого понимания.

Во-вторых, это мост между формальным и неформальным изложением. Каждое утверждение здесь
пронумеровано \emph{дословно} так же, как в оригинальной книге, и снабжено ссылкой на
соответствующее имя в Lean-файле того же раздела — мелкая серая строка \texttt{Lean:~...} под
формулировкой. По ней всегда можно свериться, как одно и то же утверждение выглядит в прозе и
в машинно проверяемом коде.

\paragraph{Что здесь есть.}
Конспект не претендует на полноту и не заменяет книгу — это личные заметки читателя. В первую
очередь фиксируются определения и доказательства, уже формализованные на Lean, и те места, где
стоило разобраться письменно, чтобы честно проверить понимание рассуждения.

\paragraph{Как это читать.}
\begin{itemize}
  \item Нумерация (\emph{Определение 1.1.1}, \emph{Упражнение 1.4.vi}, \dots) совпадает с
        нумерацией книги — по ней легко найти утверждение в оригинале.
  \item Упражнения нумеруются римскими цифрами (\emph{i}, \emph{ii}, \dots) — так устроена
        нумерация самой книги.
  \item Строка \texttt{Lean:~...} под формулировкой — это полное имя определения, леммы или
        теоремы в соответствующем файле \path{CategoryTheoryInContextLean/Section_X_Y.lean}.
  \item Курсив используется по минимуму — только там, где есть настоящее смысловое ударение.
\end{itemize}
